\section{非讼程序 - 督促程序（支付令）}
\begin{enumerate}
  \item 申请条件。
        \begin{enumerate}
          \item 请求给付的标的为金钱 \hz 有价证券。
          \item 请求给付的金钱数额确定且已经到期。
          \item 债权人没有对待给付义务。\zd{不是互负给付义务}。
          \item 必须能够送达给债务人，\zd{不能公告送达}。
          \item \zd{已经申请诉前保全的不允许申请支付令}。
          \item \zd{支付令跟担保人没有关系}。
        \end{enumerate}
  \item 向债务住所地的基层法院提出申请；符合条件 \yjt\, 15 日内发出止付令；不符合 \yjt 裁定驳回。
  \item 效力。
        \begin{enumerate}
          \item \zd{支付令作出之日起}发生法律效力 vs 产生强制执行效力不要搞混淆。
          \item 15 天内债务人清偿债务 \hz \zd{书面提出（实质）异议} \yjt 强制执行。
          \item 起诉对支付令效力的影响。
                \begin{enumerate}
                  \item 债权人起诉 \yjt 直接失效。
                  \item 债务人 \yi 向\zd{做出支付令的法院}起诉视为书面的异议 \er 向\zd{其他法院}起诉不影响支付令效力。
                \end{enumerate}
        \end{enumerate}
\end{enumerate}
